% !TEX root = ../main.tex

\section{Experiments}
\label{sec:experiments}

\subsection{Experimental Setup}
\label{sec:setup}

\noindent\textbf{Benchmarks.} We evaluate on three long-running task streams that enforce temporal order: tasks are solved using only information available before each task's release date. Table~\ref{tab:bench_stats} summarizes the statistics.

\begin{table}[t]
\centering
\small
\setlength{\tabcolsep}{4pt}
\caption{Benchmark statistics. All three enforce strict temporal order: task $i$ is solved using only information available before its release date.}
\label{tab:bench_stats}
\begin{tabular}{lrll}
\toprule
Benchmark & Tasks & Time span & Domain \\
\midrule
PolyBench & 5{,}075 & Feb 6--22, 2026 & Prediction markets \\
CTF-Dojo  & 261   & 2011--2024 & Cybersecurity CTF \\
FutureX   & 503   & Jan--Apr 2026 & Event forecasting \\
\bottomrule
\end{tabular}
\end{table}

\noindent\textbf{Solver and evolver.} The solver is Claude Sonnet 4.6 (temperature 0) and the evolver is Claude Opus 4.6 (temperature 0) across all methods. All methods share the same batch size (100 / 20 / 20 for PolyBench / CTF-Dojo / FutureX), the same batch loop, and the same temporal-reveal gate. The only axis that varies is the evolution algorithm.

\noindent\textbf{Baselines.} We compare against four systems spanning the major auto-harness paradigms:
\begin{compactenum}[$\bullet$]
\item \textbf{H0 (no evolution)}: the base solver with no harness modification.
\item \textbf{A-Evolve}~\cite{aevolve}: a linear-chain single-agent evolver that mutates prompts, skills, memory, and tools each cycle.
\item \textbf{GEPA-lite}~\cite{gepa2025}: our continual-stream port of the reflective Pareto prompt evolver, restricted to prompt-only mutations.
\item \textbf{Meta-Harness-lite}~\cite{metaharness2026}: the archive-based proposer with a growing filesystem archive and $k{=}1$ proposals per cycle.
\item \textbf{OctoTools}~\cite{octotools2026}: a human-designed Planner$+$Executor$+$Tool-Cards generalist, installed once and frozen (no self-improvement).
\end{compactenum}

\noindent\textbf{Research questions.}

\begin{compactenum}[$\bullet$]
\item \textbf{RQ1} (\S\ref{sec:main-results}): Does the full system outperform auto-harness baselines on long-running streams?
\item \textbf{RQ2} (\S\ref{sec:ablation}): Which architectural patterns of the multi-agent evolver are load-bearing for continual harness construction?
\item \textbf{RQ3} (\S\ref{sec:navloss}): Does per-task routing recover $\Ladapt$, and is the gap substantial across benchmarks?
\item \textbf{RQ4} (\S\ref{sec:hitl}): Do human-in-the-loop gains concentrate on experience-insufficient slices?
\end{compactenum}

\subsection{Main Results}
\label{sec:main-results}

Table~\ref{tab:rq1_main} reports the main comparison. Each benchmark uses metrics tailored to its practical deployment scenario:

\begin{compactenum}[(i)]
\item \emph{PolyBench} (prediction market, 5{,}075 trades over 16 days): the long-term goal is portfolio profit. We report coverage (fraction of markets traded), accuracy (correct predictions over all tasks), capital-weighted return (CWR), and portfolio return --- the metrics a trading desk would optimize.
\item \emph{CTF-Dojo} (security competition, 261 challenges over 13 years): the long-term goal is becoming a more capable hacker. We report pass rate (flags captured), CSR~$\Delta$ (whether solve rate improves or degrades as the stream progresses), and average turns per successful solve (efficiency --- a skilled hacker solves in fewer attempts).
\item \emph{FutureX} (forecasting service, 503 predictions over 82 days): the long-term goal is a reliable, broad forecasting service. We report pass rate, domain coverage (geometric mean of per-domain accuracy across 6 domains --- penalizes systems that excel on one domain but fail on others), and batch Sharpe (mean batch accuracy / std --- captures deployment reliability, since a high-stakes forecasting service needs consistent performance rather than volatile batch-to-batch swings).
\end{compactenum}

\begin{table*}[t]
\centering
\small
\setlength{\tabcolsep}{4pt}
\caption{Main results across three task streams. Bold = best per metric. Metrics are tailored to each benchmark's deployment goal (\S\ref{sec:main-results}).}
\label{tab:rq1_main}
\begin{tabular}{l@{\hskip 8pt}ccc@{\hskip 8pt}ccc@{\hskip 8pt}ccc}
\toprule
 & \multicolumn{3}{c}{\textbf{PolyBench}} & \multicolumn{3}{c}{\textbf{CTF-Dojo}} & \multicolumn{3}{c}{\textbf{FutureX}} \\
\cmidrule(lr){2-4} \cmidrule(lr){5-7} \cmidrule(lr){8-10}
Method & Coverage & Accuracy & Return & Pass & CSR$\Delta$ & Turns & Pass & DomCov & Sharpe \\
\midrule
\multicolumn{10}{l}{\emph{No evolution}} \\
Base agent      & 31.7 & 22.2 & $+$1.7 & 37.2 & $-$12.7 & 19.5 & 31.0 & 33.1 & 1.93 \\
\midrule
\multicolumn{10}{l}{\emph{Auto-harness baselines}} \\
A-Evolve        & 21.1 & 18.4 & $+$7.2 & 45.2 & $-$14.6 & \textbf{9.5} & 47.5 & 54.4 & 2.51 \\
GEPA-lite       & 32.6 & 13.4 & $+$0.3 & 42.9 & $-$11.9 & 18.5 & 28.2 & 31.7 & 1.87 \\
Meta-Harness    & 55.3 & 50.8 & $+$320 & 41.0 & $-$8.8 & 18.3 & 29.4 & 31.6 & 2.42 \\
\midrule
\multicolumn{10}{l}{\emph{Human-designed}} \\
OctoTools       & 54.6 & 39.9 & $+$19.1 & 38.3 & $-$26.5 & 17.0 & 25.6 & 28.6 & 1.77 \\
\midrule
\multicolumn{10}{l}{\emph{Our contributions}} \\
Multi-agent     & \textbf{95.8} & \textbf{79.8} & $+$351 & \textbf{52.1} & \textbf{$+$12.3} & 12.8 & \textbf{49.5} & \textbf{57.1} & 2.57 \\
Navigation      & 91.4 & 77.4 & \textbf{$+$352} & 46.0 & $-$8.8 & 15.5 & 44.1 & 50.6 & \textbf{2.73} \\
\bottomrule
\end{tabular}

\end{table*}

No single baseline dominates. A-Evolve achieves the highest baseline pass rates on CTF-Dojo and FutureX but trades only 21\% of PolyBench markets and degrades over the CTF stream (CSR~$\Delta = -14.6$). Meta-Harness-lite achieves high PolyBench coverage but trails on CTF and FutureX. Our multi-agent evolver is the \emph{only} system with positive CSR~$\Delta$ on CTF-Dojo ($+12.3$), indicating it genuinely compounds capability over the stream rather than degrading. Navigation achieves the highest batch Sharpe on FutureX (2.73), reflecting the stability that per-task routing provides: by matching each prediction to the best-suited branch, it avoids the volatile batch-to-batch swings of single-harness systems.

\subsection{Ablation: Continual Auto-Harness Internals}
\label{sec:ablation}

This subsection probes the continual-auto-harness instantiation (\S\ref{sec:method-multi}) by removing one architectural contribution at a time from the full multi-agent configuration and measuring end-of-stream performance across all three benchmarks.

Figure~\ref{fig:rq2_stats} reports the result. We compare four conditions: the single-agent A-Evolve baseline, multi-agent without temporal-reveal feedback ($-$feedback), multi-agent without persistent cross-cycle state ($-$memory), and the full system. Two findings stand out. First, persistent cross-cycle memory is the most load-bearing component: removing it causes the largest drop on PolyBench (CWR from $44.3\%$ to $7.5\%$), confirming that long-term compounding requires the evolver to build on prior cycle findings rather than restarting. Second, the full multi-agent system consistently outperforms the single-agent baseline across all benchmarks (CTF-Dojo: $43\%$ vs.\ $38\%$; FutureX: $44\%$ vs.\ $38\%$), and the gains are not explained by compute alone --- Appendix~\ref{app:compute_sanity} reports a compute-matched single-agent control that does not close the gap.

\begin{figure}[t]
\centering
\includegraphics[width=\columnwidth]{figures/figure_rq2.pdf}
\caption{\textbf{Ablation: continual harness construction.} End-of-stream performance under four ablation conditions across all three benchmarks. Bars show pass rate (CTF-Dojo, FutureX) or CWR (PolyBench) with bootstrap 95\% CI. \emph{Single-agent}: A-Evolve baseline. \emph{$-$feedback}: multi-agent without temporal-reveal gate. \emph{$-$memory}: multi-agent without persistent cross-cycle state. \emph{Full (multi-agent)}: all three contributions enabled. The full system achieves the highest performance on all benchmarks; removing persistent memory causes the largest drop on PolyBench (CWR $44.3\% \to 7.5\%$), confirming that cross-cycle state is the key driver of long-term compounding.}
\label{fig:rq2_stats}
\end{figure}


\subsection{Measuring the Adaptation Loss $\Ladapt$}
\label{sec:navloss}

To isolate the effect of per-task routing from evolver effort, we evaluate four routing strategies on the same evolved strategy tree (Figure~\ref{fig:rq3_stats}): \emph{Oracle} (per-task best branch), \emph{Adapt} (our navigator), \emph{Naive} (always \texttt{main}), and \emph{Worst} (per-task worst branch). All share identical evolver effort; only routing differs. The gap (Oracle $-$ Naive) upper-bounds $\Ladapt$.

The adaptation gap is substantial: on CTF-Dojo, Oracle achieves 55\% vs.\ Naive 18\% (37-point gap). Our navigator recovers roughly half (Adapt 35\%). On PolyBench: Oracle $+$12\% CWR vs.\ Naive $-$11\%, Adapt $+$6\%. On FutureX, Adapt (47\%) nearly matches Oracle and far exceeds Naive (22\%). The consistent ordering Oracle $>$ Adapt $>$ Naive $>$ Worst confirms that $\Ladapt$ is real and substantially recoverable by per-task routing.

\begin{figure}[t]
\centering
\includegraphics[width=\columnwidth]{figures/figure_rq3.pdf}
\caption{\textbf{Measuring $\Ladapt$: four routing strategies on the same evolved tree.} Each dot shows end-of-stream performance (pass rate or CWR) with bootstrap 95\% CI, across all three benchmarks. \emph{Oracle}: per-task best branch (upper bound on adaptation). \emph{Adapt}: our navigator routes per task. \emph{Naive}: always use \texttt{main} branch. \emph{Worst}: per-task worst branch (lower bound). All four share identical evolver effort; only the routing decision differs. The gap between Oracle and Naive upper-bounds $\Ladapt$; Adapt recovers most of it (CTF-Dojo: 35\% vs.\ Oracle 55\%; PolyBench: $+$6\% CWR vs.\ Oracle $+$12\%), confirming that per-task routing substantially reduces the adaptation loss.}
\label{fig:rq3_stats}
\end{figure}


\subsection{Human-in-the-Loop on Novel Task Regimes}
\label{sec:hitl}
